\chapter{Running External Scientific Software}
\label{ch:external-software}

Scientific workflows rarely consist only of Python functions. They typically
wrap large external programs: compiled models, command-line preprocessors,
domain-specific converters, visualization tools, and publication utilities.
DAGonStar's task model, developed across the preceding chapters, is designed
precisely for this style of file-oriented scientific computing. This chapter
collects the engineering discipline needed to wrap such software reliably,
in preparation for Part~III's environmental-science case study.

\section{External Software Prerequisites}

Before placing a program inside a DAGonStar task, it must be verified
\emph{outside} DAGonStar. For each external tool, record:

\begin{itemize}[leftmargin=1.4em]
  \item executable path;
  \item version;
  \item build/compiler information, if compiled;
  \item required environment variables;
  \item required dynamic libraries;
  \item required input files;
  \item expected output files;
  \item working-directory assumptions;
  \item license or access restrictions;
  \item expected runtime resources.
\end{itemize}

\begin{lstlisting}[style=bash, caption={Verifying a compiled executable before integration.}]
which wrf.exe
wrf.exe --help || true
ldd wrf.exe 2>/dev/null || otool -L wrf.exe 2>/dev/null
\end{lstlisting}

\begin{lstlisting}[style=bash, caption={Verifying a Python tool.}]
python -V
python -m pip freeze
python script.py --help
\end{lstlisting}

\begin{lstlisting}[style=bash, caption={Verifying a containerized tool.}]
docker run --rm image:tag command --help
\end{lstlisting}

\section{Choosing an Execution Environment}

\begin{center}
\begin{tabular}{@{}ll@{}}
\toprule
\textbf{Execution environment} & \textbf{Typical use} \\
\midrule
\texttt{TaskType.BATCH}  & Local shell command. \\
\texttt{TaskType.SLURM}  & Scheduler-submitted shell command. \\
\texttt{TaskType.DOCKER} & Command wrapped in Docker execution. \\
\texttt{TaskType.CLOUD}  & Command executed on a cloud-backed remote node. \\
\texttt{TaskType.LLM}    & OpenAI-compatible JSON request; for API-based text
  analysis rather than a locally installed command-line program. \\
\texttt{TaskType.NATIVE} & Importable Python function with scalar and staged
  file bindings. \\
\texttt{TaskType.WEB}    & Scratch-local HTTP/HTTPS request with declared
  response outputs. \\
remote variants & \texttt{BATCH}/\texttt{SLURM}/\texttt{DOCKER} selected by
  SSH parameters. \\
\bottomrule
\end{tabular}
\end{center}

\section{File-Based Data Exchange}

DAGonStar's most robust integration pattern is file exchange. Each task:

\begin{enumerate}[leftmargin=1.6em]
  \item reads files from its working directory or staged inputs;
  \item executes external software;
  \item writes named output files;
  \item is consumed by downstream tasks through \texttt{workflow://}
    references (\Cref{ch:schema}).
\end{enumerate}

This pattern is robust precisely because most scientific software already
communicates through files: NetCDF, GRIB, CSV, GeoTIFF, JSON, shapefiles,
binary restart files, and logs are the native exchange formats of
environmental and earth-system modeling, and DAGonStar requires no special
support for any of them --- only a file path.

\begin{lstlisting}[caption={A post-processing step referencing simulation output.}]
post = DagonTask(
    TaskType.BATCH,
    "postprocess",
    "python post.py --input workflow:///simulation/output/model.nc --output product.nc",
)
\end{lstlisting}

\section{Preparing Launcher Scripts}

For complex software, a long command should not be embedded directly into
the task string. Instead, create a small wrapper script and call it from
DAGonStar.

\begin{lstlisting}[style=bash, caption={\texttt{run\_model.sh}: a defensive wrapper around a compiled model.}]
#!/usr/bin/env bash
set -euo pipefail

input_file="$1"
output_dir="$2"

mkdir -p "$output_dir"

module purge 2>/dev/null || true
module load netcdf 2>/dev/null || true

export OMP_NUM_THREADS="${OMP_NUM_THREADS:-1}"

./model.exe --input "$input_file" --output "$output_dir/result.nc"

test -s "$output_dir/result.nc"
\end{lstlisting}

\begin{lstlisting}[caption={The corresponding DAGonStar task.}]
model = DagonTask(
    TaskType.BATCH,
    "model",
    "bash run_model.sh workflow:///prepare/input/forcing.nc output",
)
\end{lstlisting}

\subsection{Wrapper Script Checklist}

A good wrapper script should:

\begin{itemize}[leftmargin=1.4em]
  \item fail fast (\texttt{set -euo pipefail});
  \item print useful diagnostic information;
  \item create required output directories;
  \item validate required inputs exist;
  \item run the external program;
  \item validate expected outputs exist and are non-empty;
  \item avoid embedding secrets;
  \item write logs to predictable files.
\end{itemize}

\begin{lstlisting}[style=bash, caption={A representative input-validation block.}]
if [ ! -s "$input_file" ]; then
  echo "Missing input file: $input_file" >&2
  exit 2
fi
\end{lstlisting}

\section{Working Directories}

DAGonStar changes into the task working directory before executing the
command; external software should therefore use paths relative to the task
directory when possible. A recommended internal layout:

\begin{lstlisting}[style=bash]
task scratch/
  input/ or .dagon/inputs/
  output/
  logs/
  config/
\end{lstlisting}

\section{Environment Modules and Activation}

HPC systems commonly use environment modules:

\begin{lstlisting}[style=bash]
module purge
module load gcc/12.2 netcdf openmpi
\end{lstlisting}

Python tools often use virtual environments:

\begin{lstlisting}[style=bash]
. /path/to/venv/bin/activate
python tool.py
\end{lstlisting}

or Conda:

\begin{lstlisting}[style=bash]
source /path/to/miniconda/etc/profile.d/conda.sh
conda activate environment-name
\end{lstlisting}

\begin{practicebox}
Put this setup \emph{inside the wrapper script}, not in the DAGonStar task
command string, so that the task remains reproducible independent of the
calling shell's pre-existing environment.
\end{practicebox}

\section{Slurm and Docker Wrapper Patterns}

\begin{lstlisting}[caption={Slurm wrapper pattern: command focus stays on the model, resources are declared at construction time.}]
task = DagonTask(
    TaskType.SLURM,
    "wrf",
    "bash run_wrf.sh workflow:///real/wrfinput_d01 output",
    partition="compute",
    ntasks=64,
    time="06:00:00",
)
\end{lstlisting}

The generated \texttt{sbatch} command wraps DAGonStar's launcher script; the
wrapper script itself should still validate inputs and outputs.

\begin{lstlisting}[caption={Docker wrapper pattern.}]
task = DagonTask(
    TaskType.DOCKER,
    "map_render",
    "python map2png.py workflow:///post/product.nc output/map.png",
    image="my-lab/map-render:1.0",
)
\end{lstlisting}

The container image must contain the executable, required system libraries,
Python/R packages, fonts or geospatial data needed by renderers, and any
shell tools used by wrapper scripts.

\section{Logs and Provenance}

External software should write logs that remain inside the task scratch
directory:

\begin{lstlisting}[style=bash]
./model.exe "$input" "$output" > logs/model.stdout 2> logs/model.stderr
\end{lstlisting}

A complete provenance record additionally captures: command line; software
version; input checksums when possible; output checksums for critical
products; and the runtime environment.

\section{Error Behavior}

DAGonStar marks a task as failed when the executed command returns an error
code, or when the launcher script raises an exception. Wrapper scripts
should return non-zero exit codes when \emph{scientific} validation fails,
even if the external program itself exits with status zero:

\begin{lstlisting}[style=bash]
if ! ncdump -h output/result.nc >/dev/null; then
  echo "Invalid NetCDF output" >&2
  exit 3
fi
\end{lstlisting}

\section{The Five-Stage Practical Design Pattern}

For each external software component, the documentation recommends a
canonical five-task decomposition:

\begin{enumerate}[leftmargin=1.6em]
  \item \texttt{prepare\_*}: creates configuration and input files.
  \item \texttt{run\_*}: executes the external program.
  \item \texttt{validate\_*}: checks scientific and structural validity.
  \item \texttt{postprocess\_*}: converts outputs into analysis products.
  \item \texttt{publish\_*}: exports final artifacts.
\end{enumerate}

\texttt{workflow://} references connect each stage, so that data lineage
remains explicit end to end --- this five-stage pattern is precisely the
decomposition Part~III's environmental application case study (\Cref{ch:case-study})
will apply to a real numerical weather pipeline.

\begin{figure}[h]
\centering
\begin{tikzpicture}[
  node distance=10mm,
  task/.style={rectangle, rounded corners, draw=dagonblue, very thick,
    fill=blue!4, minimum width=24mm, minimum height=10mm, font=\sffamily\scriptsize, align=center},
  arr/.style={-{Stealth[length=2.4mm]}, thick, dagonblue}
]
  \node[task] (p) {prepare\_*};
  \node[task, right=18mm of p] (r) {run\_*};
  \node[task, right=18mm of r] (v) {validate\_*};
  \node[task, right=18mm of v] (post) {postprocess\_*};
  \node[task, right=18mm of post] (pub) {publish\_*};
  \draw[arr] (p) -- (r);
  \draw[arr] (r) -- (v);
  \draw[arr] (v) -- (post);
  \draw[arr] (post) -- (pub);
\end{tikzpicture}
\caption{The canonical five-stage wrapping pattern for external scientific
software, connected end to end by \texttt{workflow://} references.}
\label{fig:five-stage-pattern}
\end{figure}

\section{Defensive Scripting as a Research-Software-Engineering Discipline}

The wrapper-script checklist of \Cref{ch:external-software} draws on a body
of practice often called \emph{defensive programming}, and it is worth
making the underlying rationale explicit, because the checklist's individual
items can otherwise read as an arbitrary list of best practices rather than
as a coherent response to a specific category of risk.

The risk that defensive wrapper scripting addresses is the gap between
\emph{a program exiting with status zero} and \emph{a program having done
what was scientifically intended}. Most compiled scientific models,
including the WRF components discussed in \Cref{ch:case-study}, were not
originally designed with workflow-engine integration in mind, and their
exit-code conventions are frequently inherited from decades-old Fortran or C
codebases that may return a zero exit status even after writing a truncated,
corrupted, or scientifically meaningless output file, simply because the
specific failure mode encountered was never anticipated by the original
authors as something the calling environment would need to detect
programmatically. The \texttt{set -euo pipefail} preamble, input
existence validation, and output non-emptiness validation enumerated in
\Cref{ch:external-software}'s checklist are not generic shell-scripting
hygiene; each addresses a specific, named failure mode in this gap:

\begin{description}[leftmargin=2em, style=nextline]
  \item[\texttt{set -e}] ensures that a failure in any \emph{intermediate}
    command within a multi-step wrapper script propagates to the script's
    own exit status, rather than being silently masked by a later command's
    success.
  \item[\texttt{set -u}] converts a reference to an unset shell variable ---
    frequently the symptom of a missing or misnamed environment-module load
    --- from silent, wrong behavior into an immediate, loud failure.
  \item[\texttt{set -o pipefail}] ensures that a failure anywhere within a
    shell pipeline (\texttt{a | b | c}) is detected, rather than only the
    exit status of the pipeline's final command, which would otherwise mask
    an upstream failure.
  \item[Input existence/non-emptiness validation] catches the case where an
    upstream \texttt{workflow://} dependency resolved structurally
    (\Cref{ch:schema}) but the referenced file is empty or malformed ---
    a failure mode the dependency graph itself cannot detect, since it
    tracks task completion, not output content quality.
  \item[Output existence/non-emptiness validation] catches exactly the
    ``zero exit status, garbage output'' failure mode this section opened
    with, converting a silent scientific error into an explicit task
    failure DAGonStar's \texttt{Status} model (\Cref{sec:status-model}) can
    record and a downstream researcher can act on.
\end{description}

\begin{practicebox}
A useful framing for a research group adopting this discipline: every line
in the wrapper-script checklist exists to convert one specific, otherwise-
silent failure mode into an explicit, recorded \texttt{FAILED} task status.
A wrapper script that omits a checklist item has not merely skipped a
stylistic nicety; it has left one corresponding failure mode invisible to
the workflow engine, and therefore invisible to whatever automated
checkpointing, retry, or notification logic a researcher might build on top
of task status.
\end{practicebox}

\section{The Cost of Skipping External Verification}

\Cref{ch:external-software} opened with a prerequisites checklist --- record
the executable path, version, required libraries, and so on --- to be
completed \emph{before} a program is placed inside a DAGonStar task. It is
worth being explicit about why this ordering matters, since the temptation
to skip directly to workflow integration is understandable and the
consequence of doing so is not always immediately visible.

When an external program is verified manually first --- run once at a
terminal, with its output inspected by a human --- any configuration error
(a missing library, an unset environment variable, an incompatible input
format) surfaces in the simplest possible debugging context: a single
command, a single terminal, full visibility into standard output and
standard error as they are produced. When the same program is instead
integrated directly into a multi-task DAGonStar workflow without prior
manual verification, an identical configuration error instead surfaces
inside a generated launcher script, inside a scratch directory, captured to
a file rather than a terminal, and potentially obscured by whatever
upstream staging or downstream dependency behavior the workflow also
exercises. The debugging procedure of \Cref{ch:lessons-03-04} --- read
\texttt{launcher.sh}, then \texttt{context.sh}, then \texttt{stdout.txt} ---
remains available and effective in this second scenario, but it is
strictly more work than simply having ruled the configuration error out
beforehand.

\begin{practicebox}
This is the underlying rationale for the local-structural-verification
pattern already introduced for Slurm and remote tasks
(\Cref{ch:backends}), now generalized: verify the smallest possible unit ---
a single external program invocation, a single task's command, a single
task's dependency resolution --- in the simplest available context before
composing it into a larger, harder-to-debug structure. This is not specific
to DAGonStar; it is a general discipline for working with compiled
scientific software in any orchestration context, and \Cref{ch:case-study}
will apply it explicitly to the multi-stage WRF pipeline.
\end{practicebox}

\section{When the Five-Stage Pattern Should Be Adapted}

The prepare/run/validate/postprocess/publish decomposition given in
\Cref{ch:external-software} is presented as canonical, and it serves well
for the great majority of file-oriented scientific pipelines this book
addresses. It is worth being equally honest about when a researcher should
deviate from it, since presenting any pattern as universally applicable
without qualification would understate the judgment genuine workflow design
requires.

For a pipeline whose ``run'' stage is itself an ensemble of many
structurally identical simulations --- differing only in initial conditions
or parameter values, a common pattern in environmental forecasting and
uncertainty quantification --- the five-stage pattern is best applied
\emph{per ensemble member}, with an additional, ensemble-level aggregation
stage following all individual ``postprocess'' tasks, rather than attempting
to force the entire ensemble through a single ``run'' task. For a pipeline
whose validation logic is genuinely inexpensive and tightly coupled to the
run stage's own internal consistency checks, collapsing ``run'' and
``validate'' into a single task may be preferable to the traceability cost
of a separate validation task whose only content is a near-trivial check
already effectively performed by the run stage itself --- an instance of the
granularity trade-off already developed in \Cref{sec:task-granularity}.

\section{Chapter Summary}

Wrapping external scientific software in DAGonStar is, in essence, a
discipline of defensive shell scripting (fail fast, validate inputs and
outputs, log everything) combined with the engine's existing dependency and
staging mechanisms. This chapter additionally grounded each item of the
wrapper-script checklist in a specific, named failure mode it exists to
convert from silent to explicit, argued for manual external verification as
a strictly cheaper debugging context than in-workflow verification, and
discussed concrete cases --- ensemble simulations, tightly coupled run/
validate logic --- in which the canonical five-stage
prepare/run/validate/postprocess/publish pattern should be deliberately
adapted rather than applied mechanically. The five-stage pattern, and its
documented exceptions, recur throughout operational scientific pipelines and
form the organizing template Part~III now applies to a real
numerical-weather-and-ocean forecasting case study.

\begin{exercise}
Using the failure-mode framing of this chapter's first section, identify
which specific, named failure mode each of the following omissions from a
wrapper script would leave undetected: (a) omitting \texttt{set -u}; (b)
omitting output non-emptiness validation; (c) omitting \texttt{pipefail}
when the wrapper pipes the external program's output through a
post-processing filter before writing it to disk.
\end{exercise}
